\documentclass[10pt,a4paper]{amsart}
\usepackage{amsmath,amssymb,amsthm,graphicx}
\usepackage{verbatim}
\newenvironment{reference}{\comment}{\endcomment}
%\newenvironment{reference}{}{}
\newenvironment{slogan}{\comment}{\endcomment}
\newenvironment{history}{\comment}{\endcomment}
% For commutative diagrams we use Xy-pic
\usepackage[all]{xy}

% We use 2cell for 2-commutative diagrams.
\xyoption{2cell}
% Theorem environments.
%
\theoremstyle{plain}
\newtheorem{theorem}[subsection]{Theorem}
\newtheorem{proposition}[subsection]{Proposition}
\newtheorem{lemma}[subsection]{Lemma}

\theoremstyle{definition}
\newtheorem{definition}[subsection]{Definition}
\newtheorem{example}[subsection]{Example}
\newtheorem{exercise}[subsection]{Exercise}
\newtheorem{situation}[subsection]{Situation}

\theoremstyle{remark}
\newtheorem{remark}[subsection]{Remark}
\newtheorem{remarks}[subsection]{Remarks}

\numberwithin{equation}{subsection}

% Macros
%
\def\lim{\mathop{\mathrm{lim}}\nolimits}
\def\colim{\mathop{\mathrm{colim}}\nolimits}
\def\Spec{\mathop{\mathrm{Spec}}}
\def\Hom{\mathop{\mathrm{Hom}}\nolimits}
\def\Ext{\mathop{\mathrm{Ext}}\nolimits}
\def\SheafHom{\mathop{\mathcal{H}\!\mathit{om}}\nolimits}
\def\SheafExt{\mathop{\mathcal{E}\!\mathit{xt}}\nolimits}
\def\Sch{\mathit{Sch}}
\def\Mor{\mathop{\mathrm{Mor}}\nolimits}
\def\Ob{\mathop{\mathrm{Ob}}\nolimits}
\def\Sh{\mathop{\mathit{Sh}}\nolimits}
\def\NL{\mathop{N\!L}\nolimits}
\def\CH{\mathop{\mathrm{CH}}\nolimits}
\def\proetale{{pro\text{-}\acute{e}tale}}
\def\etale{{\acute{e}tale}}
\def\QCoh{\mathit{QCoh}}
\def\Ker{\mathop{\mathrm{Ker}}}
\def\Im{\mathop{\mathrm{Im}}}
\def\Coker{\mathop{\mathrm{Coker}}}
\def\Coim{\mathop{\mathrm{Coim}}}

% Boxtimes
%
\DeclareMathSymbol{\boxtimes}{\mathbin}{AMSa}{"02}

%
% Macros for moduli stacks/spaces
%
\def\QCohstack{\mathcal{QC}\!\mathit{oh}}
\def\Cohstack{\mathcal{C}\!\mathit{oh}}
\def\Spacesstack{\mathcal{S}\!\mathit{paces}}
\def\Quotfunctor{\mathrm{Quot}}
\def\Hilbfunctor{\mathrm{Hilb}}
\def\Curvesstack{\mathcal{C}\!\mathit{urves}}
\def\Polarizedstack{\mathcal{P}\!\mathit{olarized}}
\def\Complexesstack{\mathcal{C}\!\mathit{omplexes}}
% \Pic is the operator that assigns to X its picard group, usage \Pic(X)
% \Picardstack_{X/B} denotes the Picard stack of X over B
% \Picardfunctor_{X/B} denotes the Picard functor of X over B
\def\Pic{\mathop{\mathrm{Pic}}\nolimits}
\def\Picardstack{\mathcal{P}\!\mathit{ic}}
\def\Picardfunctor{\mathrm{Pic}}
\def\Deformationcategory{\mathcal{D}\!\mathit{ef}}
\usepackage{cleveref}
\crefname{originallemma}{original lemma}{original lemmas}
\Crefname{originallemma}{Original lemma}{Original lemmas}
\makeatletter
\newlabel{lemma-directed-inverse-limit-finite-sober-spectral-spaces}{{lemma-directed-inverse-limit-finite-sober-spectral-spaces}{source}{}{sourcepointer-1}{}}
\newlabel{lemma-directed-inverse-limit-finite-sober-spectral-spaces@cref}{{[originallemma][0][]lemma-directed-inverse-limit-finite-sober-spectral-spaces}{[1][0][]source}{}{}{}}
\makeatother
\crefname{originallemma}{original lemma}{original lemmas}
\Crefname{originallemma}{Original lemma}{Original lemmas}
\makeatletter
\newlabel{lemma-fibres-spectral-map-quasi-compact}{{lemma-fibres-spectral-map-quasi-compact}{source}{}{sourcepointer-2}{}}
\newlabel{lemma-fibres-spectral-map-quasi-compact@cref}{{[originallemma][0][]lemma-fibres-spectral-map-quasi-compact}{[1][0][]source}{}{}{}}
\makeatother
\crefname{originallemma}{original lemma}{original lemmas}
\Crefname{originallemma}{Original lemma}{Original lemmas}
\makeatletter
\newlabel{lemma-spectral-bijective}{{lemma-spectral-bijective}{source}{}{sourcepointer-3}{}}
\newlabel{lemma-spectral-bijective@cref}{{[originallemma][0][]lemma-spectral-bijective}{[1][0][]source}{}{}{}}
\makeatother
\crefname{originallemma}{original lemma}{original lemmas}
\Crefname{originallemma}{Original lemma}{Original lemmas}
\makeatletter
\newlabel{lemma-spectral-sub}{{lemma-spectral-sub}{source}{}{sourcepointer-4}{}}
\newlabel{lemma-spectral-sub@cref}{{[originallemma][0][]lemma-spectral-sub}{[1][0][]source}{}{}{}}
\makeatother
\title{Every spectral space embeds as a constructibly closed subset of a product of two-point spaces}
\author{The Stacks Project Authors}
\date{stacks-project:topology.tex:local-snapshot-20261012; GNU FDL 1.2 or later}
\begin{document}
\maketitle
\noindent\textit{Editorial note.} All delivered mathematical text is the author's own, copied from the pinned Stacks Project chapter source sources/stacks-topology/pinned-original-source.tex; no statement or proof step was rewritten, reordered or removed.\par\medskip
\noindent\textit{Editorial note.} Preamble provenance: the Stacks Project keeps one single shared preamble for every chapter; this card takes it directly from the staged sources/stacks-topology/preamble.tex as three verbatim spans (lines 8-12, 14-18 and 154-217), so no derived file was written and nothing was added or redefined. The wrapper class and the base packages are supplied by the assembler.\par\medskip
\noindent\textit{Editorial note.} Bibliography provenance: neither the counted statement nor its proof contains a citation (this lane rescans the whole span before writing the card), so the card carries no bibliography and no entry was reconstructed.\par\medskip
\noindent\textit{Editorial note.} Selection provenance: this card was opened under the file-based control contract (CONTROL.md), with the mechanical pre-card gate run on its statement and proof interval before the card was written.\par\medskip
\noindent\textit{Allowed standard prerequisites (editorial).}\begin{itemize}
\item The characterization of spectral spaces as inverse limits of finite sober spaces at the top of the chapter. Invoked at the start of the counted proof (source lines 4885-4890). Provides the finite quotients used to embed the space.
\item The statement that an inverse limit of directed finite sober spaces is spectral. Invoked inside the counted proof (source lines 4895-4898). Identifies the limit of the images as spectral.
\item The statement that a closed subset of a spectral space is spectral, and the comparison of limits. Invoked at the end of the counted proof (source lines 4908-4916). Compares the given space with the constructed closed subset.
\end{itemize}\medskip
\ifdefined\hypertarget\hypertarget{sourcepointer-1}{}\fi
\noindent\textit{Source reference.} Original lemma lemma-directed-inverse-limit-finite-sober-spectral-spaces: sources/stacks-topology/pinned-original-source.tex, source lines 4851--4851. Navigation pointer only: the counted text refers to the Stacks item lemma-directed-inverse-limit-finite-sober-spectral-spaces, which is an existing labeled result of this same chapter and is not part of this card. The pointer carries the target's own label as its number and its definition line as its span; its content is not reproduced here.\par\medskip
\ifdefined\hypertarget\hypertarget{sourcepointer-2}{}\fi
\noindent\textit{Source reference.} Original lemma lemma-fibres-spectral-map-quasi-compact: sources/stacks-topology/pinned-original-source.tex, source lines 4589--4589. Navigation pointer only: the counted text refers to the Stacks item lemma-fibres-spectral-map-quasi-compact, which is an existing labeled result of this same chapter and is not part of this card. The pointer carries the target's own label as its number and its definition line as its span; its content is not reproduced here.\par\medskip
\ifdefined\hypertarget\hypertarget{sourcepointer-3}{}\fi
\noindent\textit{Source reference.} Original lemma lemma-spectral-bijective: sources/stacks-topology/pinned-original-source.tex, source lines 4822--4822. Navigation pointer only: the counted text refers to the Stacks item lemma-spectral-bijective, which is an existing labeled result of this same chapter and is not part of this card. The pointer carries the target's own label as its number and its definition line as its span; its content is not reproduced here.\par\medskip
\ifdefined\hypertarget\hypertarget{sourcepointer-4}{}\fi
\noindent\textit{Source reference.} Original lemma lemma-spectral-sub: sources/stacks-topology/pinned-original-source.tex, source lines 4645--4645. Navigation pointer only: the counted text refers to the Stacks item lemma-spectral-sub, which is an existing labeled result of this same chapter and is not part of this card. The pointer carries the target's own label as its number and its definition line as its span; its content is not reproduced here.\par\medskip
\begin{lemma}
\label{lemma-spectral-closed-in-product-two-point-space}
Let $W$ be the topological space with two points, one closed,
the other not. A topological space is spectral if and only if
it is homeomorphic to a subspace of a product of
copies of $W$ which is closed in the constructible topology.
\end{lemma}

\begin{proof}
Write $W = \{0, 1\}$ where $0$ is a specialization of $1$ but not vice versa.
Let $I$ be a set. The space $\prod_{i \in I} W$ is spectral by
Lemma \ref{lemma-directed-inverse-limit-finite-sober-spectral-spaces}.
Thus we see that a subspace of $\prod_{i \in I} W$ closed in the
constructible topology is
a spectral space by Lemma \ref{lemma-spectral-sub}.

\medskip\noindent
For the converse, let $X$ be a spectral space. Let $U \subset X$ be a
quasi-compact open. Consider the continuous map
$$
f_U : X \longrightarrow W
$$
which maps every point in $U$ to $1$ and every point in $X \setminus U$ to $0$.
Taking the product of these maps we obtain a continuous map
$$
f = \prod f_U : X \longrightarrow \prod\nolimits_U W
$$
By construction the map $f : X \to Y$ is spectral. By
Lemma \ref{lemma-fibres-spectral-map-quasi-compact}
the image of $f$ is closed in the constructible topology.
If $x', x \in X$ are distinct, then since $X$ is sober either $x'$
is not a specialization of $x$ or conversely. In either case (as the
quasi-compact opens form a basis for the topology of $X$) there
exists a quasi-compact open $U \subset X$ such that $f_U(x') \not = f_U(x)$.
Thus $f$ is injective. Let $Y = f(X)$ endowed with the induced topology.
Let $y' \leadsto y$ be a specialization in $Y$ and say
$f(x') = y'$ and $f(x) = y$. Arguing as above we see that
$x' \leadsto x$, since otherwise there is a $U$ such that
$x \in U$ and $x' \not \in U$, which would imply
$f_U(x') \not \leadsto f_U(x)$.
We conclude that $f : X \to Y$ is a homeomorphism by
Lemma \ref{lemma-spectral-bijective}.
\end{proof}

\end{document}
